%% -*- coding: utf-8 -*-


\chapter{句子}

\section{基础句式及句子成分}

\subsection{主谓宾}

\begin{center}
  \begin{talltblr}[
    label = none, % Removes the "Table X:" text
    ]{
    vlines, hlines,
    % 加入支持列表环境的测量参数 (记得在导言区加载 \UseTblrLibrary{varwidth})
    % measure=vbox, stretch=-1,
    rows={halign=c, valign=m},
    columns={halign=c, valign=m},
    column{1}={0.25\linewidth, halign=c},
    column{2}={0.1\linewidth, halign=c},
    column{3}={0.15\linewidth, halign=c},
    column{4}={0.3\linewidth, halign=l},
    row{1}={font=\large\bfseries, bg=azure9, halign=c}
    }
    句子&\SetCell[c=2]{c}句子成分&&作用\\
    \SetCell[r=3]{l}Scientists conduct experiments.&主语&Scientists&句子主要说明的人或事物，也是谓语动作的执行者。\\
        &谓语&conduct&说明主语的动作，即主语“做什么”或“怎么做”。\\
        &宾语&experiments&谓语动作的承受者，即谓语动作、行为的对象。\\
  \end{talltblr}
\end{center}


\subsection{主系表}

\begin{center}
  \begin{talltblr}[
    label = none, % Removes the "Table X:" text
    ]{
    vlines, hlines,
    % 加入支持列表环境的测量参数 (记得在导言区加载 \UseTblrLibrary{varwidth})
    % measure=vbox, stretch=-1,
    rows={halign=c, valign=m},
    columns={halign=c, valign=m},
    column{1}={0.25\linewidth, halign=c},
    column{2}={0.1\linewidth, halign=c},
    column{3}={0.15\linewidth, halign=c},
    column{4}={0.3\linewidth, halign=l},
    row{1}={font=\large\bfseries, bg=azure9, halign=c}
    }
    句子&\SetCell[c=2]{c}句子成分&&作用\\
    \SetCell[r=3]{l}{The policy is a crucial measure.\\这项政策是一项关键举措。}&主语&The policy&句子主要说明的人或事物。\\
        &系动词&is&连接“主语”和“表语”，本身有一定的意义，但不能表示具体的动作。\\
        &表语&a crucial measure&描述主语的状态、性质或身份，即主语“是什么”或“怎么样”。\\
  \end{talltblr}
\end{center}


\subsection{祈使句：do/be/let\(\ldots\)}

祈使句表示请求、命令、叮嘱或者劝告等语气，通常用动词原形开头。
\begin{center}
  \begin{talltblr}[
    label = none, % Removes the "Table X:" text
    ]{
    vlines, hlines,
    % 加入支持列表环境的测量参数 (记得在导言区加载 \UseTblrLibrary{varwidth})
    % measure=vbox, stretch=-1,
    rows={halign=c, valign=m},
    columns={halign=c, valign=m},
    column{1}={0.07\linewidth, halign=c, bg=azure9},
    column{2}={0.08\linewidth, halign=l},
    column{3}={0.45\linewidth, halign=l},
    column{4}={0.25\linewidth, halign=l},
    row{1}={font=\large\bfseries, bg=azure9, halign=c}
    }
    \SetCell[c=3]{c}结构 &        &                                      & 举例            \\
    \SetCell[r=2]{c}Do型 & 肯定式 & 动词原形+（宾语）+（其他成分）       & Open the door!  \\
                         & 否定式 & Don't+动词原形+（宾语）+（其他成分） & Don't touch it. \\
    \SetCell[r=2]{c}Be型 & 肯定式 & Be+表语（名词或形容词）+（其他成分）       & Be quite!  \\
                         & 否定式 & Don't+be+表语（名词+形容词）+（其他成分） & Don't be angry with him. \\
    \SetCell[r=2]{c}Let型 & 肯定式 & Let+宾语+动词原形+（其他成分）       & Let me help you carry the box.  \\
                         & 否定式 & {Don't+let+宾语+动词原形+（其他成分）\\Let+宾语+not+动词原形+（其他成分）} & {Don't let the children play with fire.\\Let's not waste time.} \\
    \SetCell[r=2]{c}其他 & 否定式 & No+名词+V.ing, 表示“禁止做某事”       & No parking!  \\
                         & 否定式 & Never+祈使句肯定形式 & Never give up! \\
  \end{talltblr}
\end{center}


\subsection{并列句}

并列句之两个或两个以上的简单句用并列连词and、but、so、or等连接在一起后组成的句子，结构是“简单句+并列连词+简单句”，汇总如下：
\begin{center}
  \begin{talltblr}[
    label = none, % Removes the "Table X:" text
    ]{
    vlines, hlines,
    % 加入支持列表环境的测量参数 (记得在导言区加载 \UseTblrLibrary{varwidth})
    % measure=vbox, stretch=-1,
    rows={halign=c, valign=m},
    columns={halign=c, valign=m},
    % column{1}={0.15\linewidth, halign=c, bg=azure9},
    % column{2}={0.35\linewidth, halign=l},
    % column{3}={0.35\linewidth, halign=l},
    % column{4}={0.35\linewidth, halign=l},
    row{1}={font=\large\bfseries, bg=azure9, halign=c}
    }
    类型                                   & 常用并列连词                 & 用法                                           \\
    \SetCell[r=3]{c}转折、对比并列句       & but                          & “但是”                                         \\
                                           & yet                          & “但是、然而”                                   \\
                                           & while                        & “然而”                                         \\
    \SetCell[r=2]{c}因果并列句             & for                          & 表示原因，“因为”                               \\
                                           & so                           & 表示结果，“因此，所以”（so不能与becaue一起用） \\
    \SetCell[r=3]{c}选择、条件并列句       & or                           & “或者”                                         \\
                                           & neither\(\ldots\)nor\(\ldots\)         & “既不……也不……”                                 \\
                                           & either\(\ldots\)or                & “或者……或者……”                                 \\
    \SetCell[r=3]{c}并列、顺承、递进并列句 & and                          & “和；然后”                                     \\
                                           & both\(\ldots\)and                 & “不但……而且……；既……又……”                       \\
                                           & not only\(\ldots\)but (also)\(\ldots\) & “不仅……而且……”                                 \\
  \end{talltblr}
\end{center}


\subsection{主从复合句}

\begin{center}
  \begin{talltblr}[
    label = none, % Removes the "Table X:" text
    ]{
    vlines, hlines,
    % 加入支持列表环境的测量参数 (记得在导言区加载 \UseTblrLibrary{varwidth})
    % measure=vbox, stretch=-1,
    rows={halign=c, valign=m},
    columns={halign=c, valign=m},
    column{1}={0.1\linewidth, halign=c, bg=azure9},
    column{2}={0.3\linewidth, halign=l},
    column{3}={0.45\linewidth, halign=l},
    % column{4}={0.35\linewidth, halign=l},
    row{1}={font=\large\bfseries, bg=azure9, halign=c}
    }
    主从复合句&分类&举例\\
    名词性从句&{主语从句\ 宾语从句\\表语从句\ 同位语从句}&Many experts believe that AI will being profound changes to our society.\\
    定语从句&{限制性定语从句\\非限制性定语从句}&The policy which was introduced last month aims to promote economic growth.\\
    状语从句&时间、地点、原因、结果、目的、让步、条件、比较、方式\textit{状语从句}&You can find this rare plant where the climate is humid.\\
  \end{talltblr}
\end{center}

\subsection{There be句式结构}

\begin{center}
  \begin{talltblr}[
    label = none, % Removes the "Table X:" text
    ]{
    vlines, hlines,
    % 加入支持列表环境的测量参数 (记得在导言区加载 \UseTblrLibrary{varwidth})
    measure=vbox, stretch=-1,
    rows={halign=c, valign=m},
    columns={halign=c, valign=m},
    column{1}={0.2\linewidth, halign=c, bg=azure9},
    column{2}={0.65\linewidth, halign=l},
    % column{3}={0.35\linewidth, halign=l},
    % column{4}={0.35\linewidth, halign=l},
    row{1}={font=\large\bfseries, bg=azure9, halign=c}
    }
    类型&用法\\
    There be+主语+其他&{
                        \begin{enumerate}[nosep]
                          \item 通常用来表示“某处又某人或某物”，there为虚词，没有具体意义。
                          \item be为句子的谓语，be后面的名词或名词短语为句子的真正主语。
                        \end{enumerate}
                        }\\
    There be+主语+不定式&作后置定语，不定式常用主动形式表示被动意义。\\
    There+情态动词+be+主语+其他&情态动词可以是：can, may, must, should, will, 表达自己的情感。\\
    There was a time when+定语从句&意为“曾经一度，有一个时期”，用于描述个人经历。\\
  \end{talltblr}
\end{center}

\subsection{There be句式时态}

\begin{center}
  \begin{talltblr}[
    label = none, % Removes the "Table X:" text
    ]{
    vlines, hlines,
    % 加入支持列表环境的测量参数 (记得在导言区加载 \UseTblrLibrary{varwidth})
    % measure=vbox, stretch=-1,
    rows={halign=c, valign=m},
    columns={halign=c, valign=m},
    % column{1}={0.15\linewidth, halign=c, bg=azure9},
    % column{2}={0.35\linewidth, halign=l},
    % column{3}={0.35\linewidth, halign=l},
    % column{4}={0.35\linewidth, halign=l},
    row{1}={font=\large\bfseries, bg=azure9, halign=c}
    }
    时态&be动词变化形式\\
    一般现在时&There is/are…\\
    一般过去时&There was/were…\\
    一般将来时&{There will be …\\There is/are going to …}\\
    现在完成时&There has/have been …\\
    过去完成时&There had been …\\
    过去将来时&{There would be …\\There was/were going to be …}\\
  \end{talltblr}
\end{center}


\subsection{独立主格结构（特殊的状语结构）}

独立主格结构是一种特殊的状语结构，它是又名词或代词加上非谓语动词或形容词、副词、介词短语构成的。该结构有自己的“逻辑主语”，与主句的主语不同，汇总如下：
\begin{center}
  \begin{talltblr}[
    label = none, % Removes the "Table X:" text
    ]{
    vlines, hlines,
    % 加入支持列表环境的测量参数 (记得在导言区加载 \UseTblrLibrary{varwidth})
    % measure=vbox, stretch=-1,
    rows={halign=c, valign=m},
    columns={halign=c, valign=m},
    column{1}={0.25\linewidth, halign=c, bg=azure9},
    column{2}={0.55\linewidth, halign=l},
    % column{3}={0.35\linewidth, halign=l},
    % column{4}={0.35\linewidth, halign=l},
    row{1}={font=\large\bfseries, bg=azure9, halign=c}
    }
    结构&举例\\
    名词/代词+分词（现在分词/过去分词）&The task finished, they went home early.\\
    名词/动词+不定式&So many children to support, they both have to work full-time.\\
    with+名词/动词+分词/不定式&With the final exam approaching, the students are busy reviewing.\\
    There/It+being+名词/动词&It being Sunday, the library is closed.\\
  \end{talltblr}
\end{center}


\subsection{\(with\)复合结构}

\begin{enumerate}
  \item with复合结构在句中常做时间、伴随、方式、原因、条件状语等，有时还可以作后置定语，常见结构如下：
        \begin{center}
          \begin{talltblr}[
            label = none, % Removes the "Table X:" text
            ]{
            vlines, hlines,
            % 加入支持列表环境的测量参数 (记得在导言区加载 \UseTblrLibrary{varwidth})
            % measure=vbox, stretch=-1,
            rows={halign=c, valign=m},
            columns={halign=c, valign=m},
            column{1}={0.3\linewidth, halign=c, bg=azure9},
            column{2}={0.5\linewidth, halign=l},
            % column{3}={0.35\linewidth, halign=l},
            % column{4}={0.35\linewidth, halign=l},
            row{1}={font=\large\bfseries, bg=azure9, halign=c}
            }
            结构&举例\\
            with+名词/动词+介词短语&Look at the man with a book under his arm. (后置定语)\\
            with+名词/动词+形容词&The old man sat on the bench with his eyes closed. (伴随状语)\\
            with+名词/动词+副词&The plants died with the water out. (原因状语)\\
          \end{talltblr}
        \end{center}
  \item with复合结构后面还可以接动词的不同形式（to do、doing和done），用法如下：
        \begin{center}
          \begin{talltblr}[
            label = none, % Removes the "Table X:" text
            ]{
            vlines, hlines,
            % 加入支持列表环境的测量参数 (记得在导言区加载 \UseTblrLibrary{varwidth})
            % measure=vbox, stretch=-1,
            rows={halign=c, valign=m},
            columns={halign=c, valign=m},
            column{1}={0.15\linewidth, halign=c, bg=azure9},
            column{2}={0.25\linewidth, halign=l},
            column{3}={0.45\linewidth, halign=l},
            % column{4}={0.35\linewidth, halign=l},
            row{1}={font=\large\bfseries, bg=azure9, halign=c}
            }
            结构&用法&举例\\
            with+名词/动词+done&过去分词和名词/动词是被动关系，表示动作已经完成&With the problem solved, they could focus on the next task.\\
            with+名词/动词+doing&表示主动或动作正在进行&With the wind blowing strongly, the trees swayed violently.\\
            with+宾语+to do&to do和宾语是被动关系，表示尚未发生的动作&With several projects to finish this week, he has to work overtime.\\
          \end{talltblr}
        \end{center}

\end{enumerate}

\subsection{形式主语\(it\)}

It作形式主语时通常可以替代三种形式，具体如下：
\begin{enumerate}
  \item 代替不定式短语，常用于以下句型:
        \begin{itemize}
          \item “It + be +形容词+(of/for sb.) + to do sth.”,比如: It is kind of you to offer me a hand. 你能帮我一把真是太好了。
          \item “It + be +名词+to do sth,”，比如:It is your obligation to obey the law. 遵守法律是你的义务。
          \item “It takes sb. some time to do sth,”, 比如:It takes her an hour to commute to work every day. 她每天上下班通勤要花一个小时。
        \end{itemize}
  \item 代替动名词短语，常与no good、no use、useless、waste 等词连用，比如:It is useless complaining about the situation. 抱怨现状没有用。
  \item 代替主语从句，比如:It is a pity that he couldn't witness the victory. 很遗憾他没办法见证这场胜利。
\end{enumerate}

\subsection{双宾语：直接宾语和间接宾语}

宾语在句子中主要充当动作的承受者。在英语句子中，还经常会出现双宾语的情况也就是“直接宾语”和“间接宾语”，关于“双宾语”的用法汇总如下：
\begin{enumerate}
  \item 直接宾语表示动作的承受者，一般是物；间接宾语表示动作是对谁或为谁做的，一般是人。比如:I sent my friend an email.我给我的朋友发了一封电子邮件。 (其中 an emai就是直接宾语，my friend 为间接宾语)
  \item 直接宾语和间接宾语可以互换位置，比如provide sb. with sth.，互换位置后就变成provide sth. for sb.，这里互换位置需要借助介词for，还有一些谓语动词后的双宾语互换位置需要借助介词 to，如：
        \begin{itemize}
          \item book sb. sth. = book sth. for sb:为某人预定某物
          \item give sb. sth. = give sth. to sb.给某人某物
          \item lend sb. sth. = lend sth. to sb.把某物借给某人
          \item offer sb. sth. = offer sth. to sb.为某人提供某物
          \item prepare sb. sth. = prepare sth. for sb.为某人准备某物
          \item promise sb. sth = promise sth.to sb.向某人承诺某事/某物
        \end{itemize}
\end{enumerate}

\subsection{形式宾语\(it\)}

\begin{enumerate}
  \item it在作形式宾语时，通常要具备两个条件:
        \begin{enumerate}
          \item 真正的宾语是不定式、动名词或从句;
          \item 有宾语补足语。
        \end{enumerate}
  \item it作形式宾语的具体用法如下:
        \begin{enumerate}
          \item 代替不定式短语,即“make/find/think/regard\(\ldots\)+it+形容词/名词+不定式”，比如I think it essential to learn a second language.
          \item 代替动名词，比如:They regarded it no good smoking in public.
          \item 代替宾语从句，比如:We all thought it a pity that the show should have been delayed.
        \end{enumerate}
\end{enumerate}

\subsection{状语}

状语用来修饰动词、形容词、副词或整个句子，依据其表达的意义和作用的不同，主要分为下面各种状语：
\begin{center}
  \begin{longtblr}[
    label = none, % Removes the "Table X:" text
    ]{
    vlines, hlines,
    row{1}={font=\large\bfseries, bg=azure9},
    rows={halign=c, valign=m},
    columns={halign=l, valign=m}
    }
    状语     & 引导词举例                             \\
    时间状语 & when, while, as, before, after, since  \\
    地点状语   & where, in at, on                       \\
    原因状语 & because, since, as, due to             \\
    结果状语 & so/such \(\ldots\) that \(\ldots\)               \\
    目的状语 & to, so that, in order that             \\
    条件状语 & if, unless, provided that              \\
    程度状语 & 程度副词、比较级/最高级结构等          \\
    让步状语 & though, although, even though, despite \\
    方式状语 & as, as if, by                          \\
    伴随状语 & 分词短语、介词短语、形容词/名词短语    \\
    比较状语 & than, as \(\ldots\) as                      \\
  \end{longtblr}
\end{center}


\subsection{定语}

定语用来修饰名词或代词，最常见的定语类型就是“形容词定语”，如：an ancient temple。除了形容词，还有下面几类结构也可以充当定语：
\begin{center}
  \begin{talltblr}[
    label = none, % Removes the "Table X:" text
    ]{
    vlines, hlines,
    % 加入支持列表环境的测量参数 (记得在导言区加载 \UseTblrLibrary{varwidth})
    % measure=vbox, stretch=-1,
    rows={halign=c, valign=m},
    columns={halign=c, valign=m},
    column{1}={0.2\linewidth, halign=c, bg=azure9},
    column{2}={0.6\linewidth, halign=l},
    % column{3}={0.35\linewidth, halign=l},
    % column{4}={0.35\linewidth, halign=l},
    row{1}={font=\large\bfseries, bg=azure9, halign=c}
    }
    类型&举例\\
    介词短语定语&The painting in the hallway was created by a local artist.\\
    分词短语定语&The rapidly developing country faces both opportunities and challenges.\\
    不定式定语&The decision to invest in renewable energy was widely praised.\\
    定语从句&The report that highlights the importance of biodiversity was presented by the expert.\\
  \end{talltblr}
\end{center}


\subsection{同位语}

同位语常跟在名词或者代词后面，对其进行解释、说明或限定，常由名词、代词、数词或从句充当，有时也会用形容词充当，比如:
\begin{itemize}
  \item Ms. Green, an experienced journalist, wrote an in-depth report. 格林女士,一位经验丰富的记者，写了一篇深入的报道。 (名词作同位语)
  \item We ourselves are responsible for our actions. 我们自己要对我们的行为负责。 (代词作同位语)
  \item You two seem very different to me. 在我看来，你们二人有很大的不同。 (数词作同位语)
  \item Every suggestion, big or small, will be carefully considered. 每一条建议,无论重要与否，都会被认真考虑。 (形容词作同位语)
  \item The belief that hard work pays off drives him to keep going. 努力就会有回报的信念驱使他不断前进。 (从句作同位语)
\end{itemize}

\subsection{插入语}

插入语在英语中属于独立成分，起到附加解释或说明的作用，如本句中的for example。句子中常用作插入语的结构有:
\begin{enumerate}
  \item 副词或副词短语，比如 fortunately (幸运的是)、honestly (说实话)等。
  \item 形容词或形容词短语,比如hard to say (很难说)、most important of all (最重要的是)等。
  \item 介词或介词短语，比如in a word (总而言之)、in conclusion (总结)等。
  \item 不定式短语，比如to be frank (老实讲)、to tell you the truth (说实话)等。
  \item 分句，比如that is to say (也就是说)、what's more (更甚者)等。
  \item 现在分词短语，比如 generally speaking (一般来说)等。
\end{enumerate}

\subsection{直接引语和间接引语}

直接引语变间接引用时，时态变化通常遵循“时态倒退”的原则，即将直接引语中的时态向过去推移，具体变化如下：
\begin{center}
  \begin{talltblr}[
    label = none, % Removes the "Table X:" text
    ]{
    vlines, hlines,
    % 加入支持列表环境的测量参数 (记得在导言区加载 \UseTblrLibrary{varwidth})
    % measure=vbox, stretch=-1,
    rows={halign=c, valign=m},
    columns={halign=c, valign=m},
    column{1}={0.25\linewidth, halign=c, bg=azure9},
    column{2}={0.25\linewidth, halign=l},
    column{3}={0.3\linewidth, halign=l},
    % column{4}={0.35\linewidth, halign=l},
    row{1}={font=\large\bfseries, bg=azure9, halign=c}
    }
    变化&直接引语&间接引用\\
    一般现在时\(\rightarrow\)一般过去时&I am happy.&He said that he was happy.\\
    现在进行时\(\rightarrow\)过去进行时&I am reading a book.&She said that she was reading a book.\\
    一般过去时\(\rightarrow\)一般过去时或过去完成时&I saw a film yesterday.&She said that she had seen the film yesterday.\\
    过去进行时\(\rightarrow\)过去进行时或过去完成进行时&I was watching TV at that time.&He said that he was watching TV at that time.\\
    现在完成时\(\rightarrow\)过去完成时&We have finished the work.&They said that they had finished the work.\\
    过去完成时\(\rightarrow\)过去完成时&I had finished my homework before supper.&She said that she had finished her homework before supper.\\
    一般将来时\(\rightarrow\)过去将来时&I will go to the park tomorrow.&He said that he would go to the park the next day.\\
  \end{talltblr}
\end{center}
此外，直接引语变为间接引语时，还要注意人称上的变化规则：
\begin{center}
  \begin{talltblr}[
    label = none, % Removes the "Table X:" text
    ]{
    vlines, hlines,
    % 加入支持列表环境的测量参数 (记得在导言区加载 \UseTblrLibrary{varwidth})
    % measure=vbox, stretch=-1,
    rows={halign=c, valign=m},
    columns={halign=c, valign=m},
    % column{1}={0.15\linewidth, halign=c, bg=azure9},
    % column{2}={0.35\linewidth, halign=l},
    % column{3}={0.35\linewidth, halign=l},
    % column{4}={0.35\linewidth, halign=l},
    row{1}={font=\large\bfseries, bg=azure9, halign=c}
    }
    变化&特点&直接引语&间接引语\\
    \SetCell[r=3]{c}人称&\SetCell[r=3]{l}一随主，二随宾，第三人称不变化&从句主语是第一人称&从句主语和主句的人称一致\\
        &&从句主语是第二人称&从句主语和主句宾语的人称一致\\
        &&从句主语是第三人称&从句主语人称不变\\
  \end{talltblr}
\end{center}




%%% Local Variables:
%%% TeX-engine: xetex
%%% TeX-master: "../IELTS_300_Sentences"
%%% End:

% LocalWords:  vlines hlines halign valign xetex LocalWords utf sth sb
